\chapter{Tecnoloxías}
\label{chap:tecnoloxias}
\lettrine{N}este segundo capítulo explicanse as tecnoloxías utilizadas durante o proceso de desenvolvemento.

\section{Linguaxes}
\subsection*{HTML}
\acrfull{HTML} é un linguaxe de marcado estándar, utilízase na maioría das páxinas webs xa que serve para crear a estrutura destas, en conxunto con JavaScript e \acrfull{CSS} permite cambiar a apariencia e a funcionalidade.
\subsection*{JavaScript}
JavaScript é unha linguaxe de programación interpretado, permite engadir contido executable as páxinas web.
\subsection*{LaTeX}
LaTeX é un sistema de composición de textos de alta calidade. En xeral é usado para crear artículos e documentos técnicos e científicos aínda que pode utilizarse para todo tipo de documentos.
\subsection*{Python}
\subsection*{Sass}
\subsection*{SQL}
\subsection*{TypeScript}

\section{Frameworks e librarías}
\subsection*{Angular}
\subsection*{Django}
\subsection*{Ionic}
\subsection*{Spotipy}
https://github.com/plamere/spotipy/

\section{Ferramentas de desenvolvemento}
\subsection*{NeoVim}
\subsection*{Git}
\subsection*{GitLab}
\subsection*{Postman}
\subsection*{Figma}

\section{Sistema de xestión de base de datos}
\subsection*{SQLite}

\section{APIs}
\subsection*{Spotify API}
